\chapter{Color Representations, Hadrons, and Confinement}

\paragraph{A brief opening interlude.}
Before returning to group theory, allow one shameless commercial. \emph{The
Black Hole War} tells the story of the long argument with Stephen Hawking over
whether black holes destroy quantum information. It also records some of the
cruelty peculiar to theoretical physicists: not only Schrodinger's cat, but
Susskind's fish and their unfortunate trip down the toilet. With that duty
discharged, we return to Alice and Bob. Alice is still suffering, because Bob
is still making her learn group theory.

\lectureoverview{A symmetry becomes physical through its action on states. We
begin with rotations to introduce representations and generators, then repeat
the construction for color \(SU(3)\). Color representations combine into
hadrons, and their dynamics explain why isolated color is never observed.}

\section{The Group and the Things It Acts On}

There are two ideas in every symmetry discussion. First there is an abstract
group: a collection of operations with a multiplication law, an identity, and
inverses. Second there are the objects on which those operations act. The two
ideas can be distinguished, but in physics they cannot be separated. A
symmetry matters because it transforms physical states, fields, coordinates,
or observables.

\begin{classroomqa}
\audiencequestion{A group has been described both as an abstract multiplication table and as matrices acting on physical objects. Can those two parts be dissociated?
\lecturetimestamp{00:01:58}}
\lecturerresponse{They are different levels of description, but they are completely related. A rotation exists as an abstract operation before we specify what it rotates. Physics then asks how that same operation acts on the states and quantities relevant to a particular system. The matrices supply that concrete action while preserving the abstract multiplication law.}
\end{classroomqa}

Spatial rotations provide the cleanest example. Classically, the same rotation
may act on a particle's coordinates, on the direction of a magnet, or on any
other vector. In quantum mechanics, states live in a linear vector space, so a
rotation is represented by a linear operator on that space. If \(g\) denotes
an abstract group element and \(D(g)\) its matrix representative, consistency
means
\begin{equation}
D(g_1g_2)=D(g_1)D(g_2),
\qquad
D(e)=\boldsymbol 1.
\end{equation}
The matrices need not look like the abstract operations, but their
multiplication must reproduce the group multiplication.
\editorialnote{The homomorphism equation writes explicitly the correspondence between group multiplication and matrix multiplication described verbally in the source.}

The dimension of the matrix is determined by the state space, not by the
abstract group alone. If we retain only the spin of one electron, there are two
orthogonal possibilities, spin up and spin down, and a state is a two-component
column vector. Rotations therefore act through \(2\times2\) matrices. For two
spin-one-half objects there are four basis states,
\begin{equation}
\ket{\uparrow\uparrow},\qquad
\ket{\uparrow\downarrow},\qquad
\ket{\downarrow\uparrow},\qquad
\ket{\downarrow\downarrow},
\end{equation}
and the induced action on the combined space is four-dimensional.

The familiar sequence is
\begin{equation}
\begin{array}{c|c|c}
\text{spin} & \text{number of states} & \text{matrix dimension}\\
\hline
0 & 1 & 1\times1\\
\frac12 & 2 & 2\times2\\
1 & 3 & 3\times3\\
\frac32 & 4 & 4\times4
\end{array}
\end{equation}
and so forth. A spin-zero state furnishes the trivial representation: every
rotation leaves it alone. The one-by-one matrix is simply \(1\). A spin-one
state has three components, and its rotation matrices are the familiar
three-dimensional vector rotations. More generally, spin \(j\) has \(2j+1\)
states.

The representation can also be infinite-dimensional. If we describe the full
orbital motion of an electron rather than only its spin, a state contains an
amplitude for every position, or equivalently for every momentum. The rotation
operator then acts on an infinite-dimensional space of wavefunctions. Thus one
and the same abstract rotation group has trivial, two-dimensional,
three-dimensional, and infinite-dimensional representations. This is the fact
we will reuse when the discussion moves from spin to color.

\section{Why Spin Rotations Are Special Unitary}

Return now to a single electron spin and discard its orbital motion. A rotation
is represented by a \(2\times2\) unitary matrix \(U\),
\begin{equation}
U^\dagger U=\boldsymbol 1,
\qquad
U^{-1}=U^\dagger.
\label{eq:l4-unitarity}
\end{equation}
The inverse transformation must undo the rotation, and unitarity preserves the
norm, hence every probability.

A general \(2\times2\) complex matrix contains four complex entries, or eight
real parameters. Equation~\eqref{eq:l4-unitarity} imposes four independent real
conditions, leaving four parameters. An ordinary spatial rotation has only
three: two parameters choose the axis and one chooses the angle. The unitary
matrix therefore contains one phase too many. The additional restriction is
\begin{equation}
\det U=1.
\label{eq:l4-det-one}
\end{equation}
The count is then
\begin{equation}
8-4-1=3.
\end{equation}

The determinant-one restriction is stable under multiplication because
\begin{equation}
\det(AB)=\det A\,\det B.
\end{equation}
If two matrices each have determinant one, so does their product. It is also
compatible with unitarity. Taking the determinant of
Equation~\eqref{eq:l4-unitarity} gives
\begin{equation}
1=\det(U^\dagger U)
=\det(U^\dagger)\det U
=(\det U)^*\det U.
\end{equation}
Consequently the determinant of a unitary matrix lies on the unit circle,
\begin{equation}
\det U=e^{i\phi}.
\end{equation}

\begin{classroomqa}
\audiencequestion{Does unitarity require the determinant to be real?
\lecturetimestamp{00:12:10}}
\lecturerresponse{No. Unitarity requires only unit modulus. The determinant may be any pure phase \(e^{i\phi}\). For a two-dimensional matrix, multiplying \(U\) by \(e^{-i\phi/2}\) removes that common phase and produces a determinant-one matrix.}
\end{classroomqa}

Indeed, if
\begin{equation}
V=e^{-i\phi/2}U,
\end{equation}
then \(V\) remains unitary and
\begin{equation}
\det V=e^{-i\phi}\det U=1.
\end{equation}
The determinant-one unitary matrices form \(SU(2)\).

There is one fine point. It is tempting to say that \(SU(2)\) is simply the
group of rotations in three-dimensional space, but the correspondence is
two-to-one. The matrices \(U\) and \(-U\) act differently on a spinor, yet they
correspond to the same rotation of an ordinary vector. Thus
\begin{equation}
SU(2)/\{\boldsymbol 1,-\boldsymbol 1\}\simeq SO(3).
\end{equation}
For the present discussion the local multiplication structure and the three
generators are what matter, but the double cover explains the two-valuedness of
spin-one-half rotations.

\section{Infinitesimal Transformations and Their Generators}

A finite rotation can be built from many very small rotations. A group element
close to the identity is called infinitesimal, and knowing all infinitesimal
directions is enough, in a connected group, to reconstruct finite
transformations by repeated multiplication. More importantly for physics, the
operators that generate these infinitesimal transformations are the conserved
quantities associated with the symmetry.

Write a near-identity element as
\begin{equation}
U=\boldsymbol 1+i\epsilon T,
\qquad |\epsilon|\ll1.
\label{eq:l4-infinitesimal-u}
\end{equation}
The factor of \(i\) is a convenient convention. Substituting into unitarity and
discarding terms of order \(\epsilon^2\) gives
\begin{align}
(\boldsymbol 1+i\epsilon T)
(\boldsymbol 1-i\epsilon T^\dagger)&=\boldsymbol 1,
\\
i\epsilon(T-T^\dagger)&=0,
\\
T&=T^\dagger.
\label{eq:l4-hermitian-generator}
\end{align}
The generator is Hermitian, as a quantum-mechanical observable must be.

\begin{figure}[htbp]
\centering
\includegraphics[width=0.62\textwidth]{lecture_04_figure_03.png}
\caption{The near-identity unitarity calculation ending in \(T=T^\dagger\).}
\lectureframe{00:23:28}
\end{figure}

For spatial rotations there are three linearly independent generators because
there are three independent infinitesimal axes. A small rotation about an
arbitrary direction is assembled from small rotations about \(x\), \(y\), and
\(z\). The corresponding generators are the angular momentum operators
\(J_x,J_y,J_z\). Their commutators encode the noncommutativity of rotations,
\begin{equation}
[J_i,J_j]=i\epsilon_{ijk}J_k.
\label{eq:l4-angular-momentum-algebra}
\end{equation}
Repeated infinitesimal rotations then build a finite one.

We have used unitarity but have not yet used determinant one. For any small
matrix \(m\),
\begin{equation}
\det(\boldsymbol 1+\epsilon m)
=1+\epsilon\,\operatorname{tr}m+O(\epsilon^2).
\label{eq:l4-det-expansion}
\end{equation}

\begin{figure}[htbp]
\centering
\includegraphics[width=0.62\textwidth]{lecture_04_figure_02.png}
\caption{The first-order determinant expansion as it begins on the board.}
\lectureframe{00:21:51}
\end{figure}

Applying Equation~\eqref{eq:l4-det-expansion} to
Equation~\eqref{eq:l4-infinitesimal-u},
\begin{equation}
\det U
=1+i\epsilon\,\operatorname{tr}T+O(\epsilon^2).
\end{equation}
The special-unitary condition therefore requires
\begin{equation}
\operatorname{tr}T=0.
\label{eq:l4-traceless-generator}
\end{equation}
The generators of \(SU(2)\) are precisely the traceless Hermitian
\(2\times2\) matrices. That vector space has dimension three. A standard basis
is supplied by the Pauli matrices,
\begin{equation}
\begin{aligned}
\sigma_x&=
\begin{pmatrix}0&1\\1&0\end{pmatrix},
&
\sigma_y&=
\begin{pmatrix}0&-i\\i&0\end{pmatrix},
\\[0.4em]
\sigma_z&=
\begin{pmatrix}1&0\\0&-1\end{pmatrix}.
\end{aligned}
\end{equation}
For a spin-one-half representation one may take \(J_i=\sigma_i/2\), so a
finite rotation through angle \(\theta\) around \(\hat{\boldsymbol n}\) is
\begin{equation}
U(\hat{\boldsymbol n},\theta)
=\exp\!\left(-i\theta\,
\hat{\boldsymbol n}\cdot\boldsymbol\sigma/2\right).
\end{equation}
\editorialnote{The finite exponential is the standard completion of the infinitesimal-generator discussion. The source emphasizes repeated small rotations rather than writing this formula explicitly.}

\begin{classroomqa}
\audiencequestion{So there are three independent traceless Hermitian \(2\times2\) matrices?
\lecturetimestamp{00:24:07}}
\lecturerresponse{Yes. They are the three Pauli matrices. For \(3\times3\) special-unitary matrices the corresponding count is eight, and a conventional basis is given by the eight Gell-Mann matrices.}
\end{classroomqa}

\begin{classroomqa}
\audiencequestion{Is that eight the same ``eight'' as the eightfold way and the eight gluons?
\lecturetimestamp{00:24:43}}
\lecturerresponse{The mathematics is the same kind of \(SU(3)\) representation theory, but one must specify which \(SU(3)\) symmetry is acting. The historical eightfold-way flavor symmetry is not color \(SU(3)\). The eight gluons belong to the adjoint representation of color \(SU(3)\).}
\end{classroomqa}

\section{From SU(2) to Color SU(3)}

Now increase the matrices from \(2\times2\) to \(3\times3\). A general complex
\(3\times3\) matrix has nine complex entries, hence eighteen real parameters.
The condition \(U^\dagger U=\boldsymbol 1\) supplies nine real constraints,
and \(\det U=1\) removes one more. Therefore
\begin{equation}
18-9-1=8.
\end{equation}
There are eight independent infinitesimal directions and eight traceless
Hermitian generators. Unlike \(SU(2)\), this group is not the rotation group
of ordinary three-dimensional space. It is an internal symmetry acting on
color.

Before assigning those representations to quarks, recall how composite
representations arise. Put two spin-one-half systems on top of one another and
ignore their orbital motion. There are four product states, but a rotation does
not mix all four irreducibly. Three combinations transform as spin one,
\begin{align}
\ket{1,1}&=\ket{\uparrow\uparrow},\\
\ket{1,0}&=\frac{1}{\sqrt2}
(\ket{\uparrow\downarrow}+\ket{\downarrow\uparrow}),\\
\ket{1,-1}&=\ket{\downarrow\downarrow},
\end{align}
while the fourth transforms as spin zero,
\begin{equation}
\ket{0,0}=\frac{1}{\sqrt2}
(\ket{\uparrow\downarrow}-\ket{\downarrow\uparrow}).
\end{equation}
Thus
\begin{equation}
\boldsymbol 2\otimes\boldsymbol 2
=\boldsymbol 3\oplus\boldsymbol 1.
\label{eq:l4-two-spin-decomposition}
\end{equation}
The aligned states give the \(m=\pm1\) members of the triplet, the symmetric
zero-projection state completes it, and the antisymmetric combination is the
singlet. This is the general logic behind adding angular momenta and behind
combining color representations.

The same singlet--triplet distinction occurs whenever two spin-one-half
degrees of freedom are coupled. In atomic hydrogen the electron and proton
spins form hyperfine singlet and triplet states. Ortho- and para-hydrogen use
the same representation theory for the nuclear spins of the two protons in an
\(H_2\) molecule.
\editorialnote{The source uses ``ortho and para-hydrogen'' while referring to proton--electron spins in a hydrogen atom. The mathematical decomposition is correct; the molecular and atomic names are distinguished here for precision.}

\subsection{The Triplet and the Anti-Triplet}

A single quark has three color amplitudes. In a chosen basis write
\begin{equation}
q=
\begin{pmatrix}
q_r\\q_g\\q_b
\end{pmatrix},
\qquad
q'=Uq,
\qquad U\in SU(3).
\label{eq:l4-quark-triplet}
\end{equation}
The three entries may be state amplitudes for one quark or three field
operators that create red, green, and blue quarks. The group mixes the three
components. A state that is initially red generally becomes a superposition
of red, green, and blue after a color transformation. This fundamental
three-dimensional representation is called the \(\boldsymbol 3\).

For an anti-quark, the group acts through the complex conjugate of the quark
representation. Taking the complex conjugate of
Equation~\eqref{eq:l4-quark-triplet} gives
\begin{equation}
\bar q'=U^*\bar q.
\end{equation}
This representation is denoted \(\bar{\boldsymbol 3}\). The \(U(1)\) analogy is
immediate: if an electron field acquires \(e^{i\theta}\), a positron field
acquires \(e^{-i\theta}\). Quark and anti-quark color transformations are the
non-Abelian version of that conjugate pairing.

Near the identity, if
\begin{equation}
U=\boldsymbol 1+i\epsilon T,
\end{equation}
then
\begin{equation}
U^*=\boldsymbol 1-i\epsilon T^*.
\end{equation}
Thus the anti-triplet generators are \(-T^*\), the conjugates with the opposite
sign. This is the precise sense in which anti-quarks carry opposite color
charges.

\begin{classroomqa}
\audiencequestion{What happens to the complex-conjugate representation in the \(SU(2)\) case?
\lecturetimestamp{00:36:19}}
\lecturerresponse{The \(SU(2)\) doublet is a special case: its complex conjugate is equivalent to the original doublet after a fixed change of basis involving the antisymmetric epsilon tensor. For \(SU(3)\), the \(\boldsymbol 3\) and \(\bar{\boldsymbol 3}\) are genuinely inequivalent representations.}
\end{classroomqa}
\editorialnote{The mathematically precise term for the \(SU(2)\) fundamental representation is pseudoreal, not real. The response preserves the source's intended distinction while correcting the terminology.}

\section{Combining Color Representations}

The symbol \(\otimes\) does not mean ordinary multiplication of the numbers
written beside it. It means that two state spaces are combined. The resulting
tensor-product space is then reorganized into pieces that do not mix with one
another under the group action.

\subsection{Two Quarks: \texorpdfstring{\(\boldsymbol 3\otimes\boldsymbol 3\)}{3 x 3}}

Two quarks have nine color basis states: red--red, red--green, red--blue, and
the six remaining ordered pairs. Organize the two color indices into a
symmetric and an antisymmetric part. The symmetric part has six independent
components; the antisymmetric part has three. Under \(SU(3)\),
\begin{equation}
\boldsymbol 3\otimes\boldsymbol 3
=\boldsymbol 6\oplus\bar{\boldsymbol 3}.
\label{eq:l4-three-times-three}
\end{equation}

\begin{figure}[htbp]
\centering
\includegraphics[width=0.68\textwidth]{lecture_04_figure_04.png}
\caption{The two-quark product split into an anti-triplet and a sextet.}
\lectureframe{00:39:33}
\end{figure}

The anti-triplet may be represented explicitly by contracting two quark
indices with the invariant epsilon tensor,
\begin{equation}
(qq)^k_{\bar 3}
=\frac{1}{\sqrt2}\epsilon^{kij}q_iq_j.
\end{equation}
This object has three components and transforms like an anti-quark. The six
symmetric combinations form the sextet. The statement is not that two quarks
become an anti-quark dynamically; it is that one sector of their color wave
function has the same transformation law as an anti-quark.

\subsection{A Quark and an Anti-Quark: \texorpdfstring{\(\boldsymbol 3\otimes\bar{\boldsymbol 3}\)}{3 x anti-3}}

A quark--anti-quark color state can be viewed as a \(3\times3\) matrix
\(M_i{}^j=q_i\bar q^j\). Every such matrix splits uniquely into its trace and
traceless parts,
\begin{equation}
M=\frac{\operatorname{tr}M}{3}\boldsymbol 1
+\left(M-\frac{\operatorname{tr}M}{3}\boldsymbol 1\right).
\end{equation}
The trace is a single invariant component and the traceless part has eight
components. Therefore
\begin{equation}
\boldsymbol 3\otimes\bar{\boldsymbol 3}
=\boldsymbol 1\oplus\boldsymbol 8.
\label{eq:l4-three-times-threebar}
\end{equation}
Up to normalization, the singlet is
\begin{equation}
\ket{1}_{q\bar q}
=\frac{1}{\sqrt3}
(\ket{r\bar r}+\ket{g\bar g}+\ket{b\bar b}).
\end{equation}
It cannot rotate into anything else because it is invariant. The other eight
linear combinations mix among themselves. They furnish the adjoint
representation, the same representation in which the eight generators
transform.

\subsection{Three Quarks}

Combine a third quark with Equation~\eqref{eq:l4-three-times-three}. The
anti-triplet branch gives
\begin{equation}
\bar{\boldsymbol 3}\otimes\boldsymbol 3
=\boldsymbol 1\oplus\boldsymbol 8,
\end{equation}
so the three-quark product contains a singlet. In full,
\begin{equation}
\boldsymbol 3\otimes\boldsymbol 3\otimes\boldsymbol 3
=\boldsymbol 1\oplus\boldsymbol 8\oplus\boldsymbol 8
\oplus\boldsymbol{10}.
\label{eq:l4-three-quark-full}
\end{equation}
The color singlet is the fully antisymmetric red--green--blue combination,
\begin{equation}
\ket{1}_{qqq}
=\frac{1}{\sqrt6}\epsilon^{ijk}\ket{q_iq_jq_k}.
\label{eq:l4-baryon-singlet}
\end{equation}
Because \(\epsilon^{ijk}\) vanishes whenever two indices coincide, every
nonzero term contains one red, one green, and one blue quark, with alternating
signs under interchange.

\begin{classroomqa}
\audiencequestion{Are the non-singlet combinations merely unlikely, so that their probabilities can be ignored?
\lecturetimestamp{00:46:02}}
\lecturerresponse{No. The issue is energy, not a small probability. An isolated non-singlet would carry a color field to arbitrarily large distance and, in the confining regime, would require unbounded energy. The dynamics behind that statement will be developed below.}
\end{classroomqa}

\subsection{The Gluon Octet}

The gluon field transforms like the octet in
Equation~\eqref{eq:l4-three-times-threebar}. With respect to color it carries
one triplet index and one anti-triplet index, with the singlet trace removed.
Schematically,
\begin{equation}
q\sim\boldsymbol 3,
\qquad
\bar q\sim\bar{\boldsymbol 3},
\qquad
G\sim\boldsymbol 8.
\end{equation}
These symbols name transformation laws; particles are not literally equal to
the integers three and eight.

The ninth quark--anti-quark color combination, proportional to the identity,
is a singlet and does not mix with the octet. Omitting it is therefore
mathematically consistent: an \(SU(3)\) gauge field is traceless and has eight
components rather than the nine of a \(U(3)\) gauge field.

\begin{classroomqa}
\audiencequestion{Why does nature omit the extra singlet gluon?
\lecturetimestamp{00:49:00}}
\lecturerresponse{The immediate group-theory answer is that the singlet is a separate invariant sector, so the octet closes consistently without it. The deeper reason for choosing \(SU(3)\) rather than \(U(3)\) is not derived here; it is part of the gauge structure inferred from QCD phenomenology.}
\end{classroomqa}

\section{Color Singlets as Hadrons}

The free, finite-mass particles of the strong interaction are color singlets.
This does not make quarks or gluons unreal; it means they do not occur as
isolated asymptotic particles. The elementary singlet-building patterns are
\begin{equation}
qqq,
\qquad
q\bar q,
\qquad
GG\ \text{and higher gluonic products}.
\end{equation}
They lead respectively to baryons, mesons, and glueballs. Products of several
already colorless objects are of course colorless again, so more complicated
multihadron states can be built from these ingredients.

\begin{classroomqa}
\audiencequestion{What are the three-quark and quark--anti-quark singlets called, and can gluons also make singlets?
\lecturetimestamp{00:51:38}}
\lecturerresponse{A three-quark singlet is a baryon; a quark--anti-quark singlet is a meson. Gluons can also combine into color singlets. The simplest such strongly interacting states are called glueballs, and more than two gluons may participate.}
\end{classroomqa}

Three spin-one-half quarks always combine to an overall half-integer spin,
although orbital angular momentum can change which half-integer occurs.
Quark--anti-quark mesons have integer spin. Purely gluonic glueballs also have
integer spin.

For two gluons the first board statement is only a count,
\begin{equation}
8\times8=1+63.
\end{equation}
There are sixty-four color states in all; one combination is a singlet and the
remaining sixty-three belong to nonsinglet sectors.

\begin{figure}[htbp]
\centering
\includegraphics[width=0.68\textwidth]{lecture_04_figure_05.png}
\caption{The board's two-gluon state count: one singlet among sixty-four states.}
\lectureframe{00:53:44}
\end{figure}

The complete irreducible decomposition is
\begin{equation}
\boldsymbol 8\otimes\boldsymbol 8
=\boldsymbol 1\oplus\boldsymbol 8_S\oplus\boldsymbol 8_A
\oplus\boldsymbol{10}\oplus\overline{\boldsymbol{10}}
\oplus\boldsymbol{27}.
\label{eq:l4-eight-times-eight}
\end{equation}
The dimensions add correctly:
\begin{equation}
1+8+8+10+10+27=64.
\end{equation}
The two octets are distinct invariant subspaces, one symmetric and one
antisymmetric under exchange of the two adjoint factors. The singlet is the
two-gluon color channel from which a glueball state can be formed.
\editorialnote{The board first gives only \(8\times8=1+63\). The later classroom exchange supplies the listed irreducible dimensions; bars and symmetry labels are restored here in standard \(SU(3)\) notation.}

\subsection{Why Observable Electric Charges Are Integral}

The quark charges come in two types,
\begin{equation}
Q_{\mathrm{up\text{-}like}}=+\frac23,
\qquad
Q_{\mathrm{down\text{-}like}}=-\frac13.
\end{equation}
Three quarks therefore have an integral total charge: \(-1,0,+1,\text{ or }+2\),
depending on how many are up-like. A quark and an anti-quark likewise have
charge \(0\) or \(\pm1\). Gluons carry no electric charge. Thus the basic color
singlets built from the observed quark representations have integer electric
charge, even though their constituents carry fractional charge.

\begin{classroomqa}
\audiencequestion{What is the electric charge of the strange quark?
\lecturetimestamp{00:56:08}}
\lecturerresponse{It is down-like, so its charge is \(-1/3\). The three families repeat the same charge pattern: \((u,d)\), \((c,s)\), and \((t,b)\), with charges \(+2/3\) and \(-1/3\) in each pair.}
\end{classroomqa}

Color singletness and integral electric charge are closely linked for this
particle content, but they are not identical mathematical statements. The
integrality also uses the particular Standard Model electric-charge
assignments just written.

\begin{classroomqa}
\audiencequestion{Is a color singlet an entangled state, and does that have the same significance as an entangled pair in a Bell experiment?
\lecturetimestamp{00:56:58}}
\lecturerresponse{The singlet color wavefunction is certainly entangled: it is a coherent superposition, not a product of separately definite colors. In principle its constituents may be separated, but confinement makes that separation increasingly expensive and eventually produces new hadrons. It is therefore not a convenient system for a Bell-style manipulation, even though the state is entangled.}
\end{classroomqa}

\section{What the Tensor-Product Equations Mean}

At this point the multiplication tables deserve a more concrete explanation.
Let \(\psi_i\) and \(\phi_j\) be operators that create two spin-one-half
objects; they may describe two species, such as an electron and a proton, or
two copies of the same species when the appropriate statistics are included.
Each index takes two values. The four products
\begin{equation}
\psi_\uparrow\phi_\uparrow,
\quad
\psi_\uparrow\phi_\downarrow,
\quad
\psi_\downarrow\phi_\uparrow,
\quad
\psi_\downarrow\phi_\downarrow
\end{equation}
span the tensor-product space.

If a group element acts as \(U\) on each factor, it acts as \(U\otimes U\) on
the product. In components,
\begin{equation}
\psi_i\phi_j
\longmapsto
U_i{}^kU_j{}^\ell\psi_k\phi_\ell.
\label{eq:l4-product-action}
\end{equation}
This is a four-dimensional representation, but it is reducible. The
antisymmetric combination
\begin{equation}
\frac{1}{\sqrt2}
(\psi_\uparrow\phi_\downarrow
-\psi_\downarrow\phi_\uparrow)
\end{equation}
maps into itself, while the three symmetric combinations map among
themselves. In the basis adapted to those combinations, every product
representation matrix takes block-diagonal form,
\begin{equation}
D_{2\otimes2}(g)
\sim
\begin{pmatrix}
1&0\\
0&D_3(g)
\end{pmatrix}.
\label{eq:l4-spin-block-diagonal}
\end{equation}

\begin{classroomqa}
\audiencequestion{What operation is represented by an equation such as \(\boldsymbol 3\otimes\boldsymbol 3=\bar{\boldsymbol 3}\oplus\boldsymbol 6\)?
\lecturetimestamp{00:57:51}}
\lecturerresponse{The left-hand side forms the composite state space. The group acts on both factors and therefore mixes the product basis states. The right-hand side says that a change of basis separates those states into invariant sets: three antisymmetric combinations that transform as an anti-triplet and six symmetric combinations that transform as a sextet. The plus sign means a direct sum of alternatives, not ordinary numerical addition.}
\end{classroomqa}

\begin{classroomqa}
\audiencequestion{How do we know the nine states split as \(3+6\), or \(1+8\), rather than, say, \(2+7\)?
\lecturetimestamp{01:04:39}}
\lecturerresponse{Dimension counting alone cannot decide the split. One must know which irreducible representations \(SU(3)\) possesses and construct invariant subspaces, for example by symmetrizing indices, antisymmetrizing them, or separating trace from traceless parts. A full classification takes more group theory; here we use only the representations needed for quarks and gluons.}
\end{classroomqa}

The invariant epsilon tensor makes the singlet construction transparent. For
\(SU(2)\),
\begin{equation}
\epsilon_{ij}U^i{}_kU^j{}_\ell
=(\det U)\epsilon_{k\ell}
=\epsilon_{k\ell}.
\end{equation}
Therefore \(\epsilon_{ij}\psi^i\phi^j\) is invariant. For \(SU(3)\), the
three-index version gives
\begin{equation}
\epsilon_{ijk}U^i{}_aU^j{}_bU^k{}_c
=(\det U)\epsilon_{abc}
=\epsilon_{abc},
\end{equation}
which proves the invariance of the three-quark state in
Equation~\eqref{eq:l4-baryon-singlet}.

\begin{classroomqa}
\audiencequestion{Can one show explicitly why the singlet does not transform?
\lecturetimestamp{01:07:18}}
\lecturerresponse{Yes. Contract all fundamental indices with the fully antisymmetric epsilon tensor. Under the transformation, the product of matrices multiplies epsilon by \(\det U\). For a special-unitary matrix that determinant is one, so the contracted state is unchanged. In \(SU(2)\) this proves the two-spin singlet; in \(SU(3)\) it proves the red--green--blue three-quark singlet.}
\end{classroomqa}

\begin{classroomqa}
\audiencequestion{What is the full decomposition of two gluon octets?
\lecturetimestamp{01:09:24}}
\lecturerresponse{It is the decomposition in Equation~\eqref{eq:l4-eight-times-eight}: one singlet, two octets, a decuplet, an anti-decuplet, and a twenty-seven-dimensional representation. There is exactly one color-singlet combination of the two octet indices.}
\end{classroomqa}

\begin{classroomqa}
\audiencequestion{What is the corresponding full decomposition for three quarks?
\lecturetimestamp{01:11:51}}
\lecturerresponse{The physically essential answer is that it contains one singlet. Completing the arithmetic gives Equation~\eqref{eq:l4-three-quark-full}: one singlet, two octets, and one decuplet, for a total of \(1+8+8+10=27\) states.\editorialnote{The lecturer does not recall the complete \(3\otimes3\otimes3\) decomposition during the exchange. The standard completion is supplied explicitly and is labeled here rather than attributed to the spoken answer.}}
\end{classroomqa}

The same idea scales to the two-gluon space. In the naive product basis,
\(SU(3)\) acts through \(64\times64\) matrices. In a basis adapted to the
irreducible components, those matrices are block diagonal,
\begin{equation}
\begin{aligned}
D_{8\otimes8}(g)
&\sim
\operatorname{diag}\bigl(
1,D_{8_S}(g),D_{8_A}(g),
\\
&\hspace{4.8em}
D_{10}(g),D_{\overline{10}}(g),D_{27}(g)
\bigr).
\end{aligned}
\end{equation}
No block mixes with another, although states within each block mix among
themselves.

\begin{classroomqa}
\audiencequestion{Is tensor-product decomposition just a change to a basis in which the large representation matrix becomes block diagonal?
\lecturetimestamp{01:12:02}}
\lecturerresponse{Exactly. For two spins the four-dimensional matrix becomes a one-dimensional singlet block and a three-dimensional triplet block. For two gluons the sixty-four-dimensional matrix becomes blocks of dimensions \(1,8,8,10,10,\text{ and }27\). Each block is closed under every group transformation.}
\end{classroomqa}

\begin{classroomqa}
\audiencequestion{Is there a name for one of those closed pieces inside the full tensor-product space?
\lecturetimestamp{01:14:55}}
\lecturerresponse{It is an invariant subspace carrying a subrepresentation. When it cannot be decomposed further, it carries an irreducible representation, or irrep. The full tensor product is the direct sum of those irreducible components.\editorialnote{The source leaves the requested term unresolved in the room. Standard representation-theory terminology is supplied in the response.}}
\end{classroomqa}

\begin{classroomqa}
\audiencequestion{When the nonsinglet tens and other representations are called nonexistent, does that mean they should be ignored altogether?
\lecturetimestamp{01:15:36}}
\lecturerresponse{They are legitimate color sectors and matter in intermediate calculations, but they do not appear as isolated finite-energy particles in the laboratory. Confinement, not a failure of representation theory, prevents them from becoming free asymptotic states.}
\end{classroomqa}

\section{From Nambu's Color Idea to Confinement}

The color idea predates the final formulation of QCD. Yoichiro Nambu already
understood in the early 1960s that a hidden three-valued quantum number could
solve a striking problem: quarks carry fractional electric charge, yet no free
fractionally charged particles are observed. If physical states are forced to
be color singlets, quarks occur in the integral-charge combinations described
above. The remaining question is dynamical: what forces all finite-energy free
states into singlets?

Begin with the electromagnetic analogy. A charge and an opposite charge
attract. Bringing them together lowers the energy; pushing like charges
together requires work. The field itself stores positive energy,
\begin{equation}
E_{\mathrm{field}}
=\frac12\int d^3x\,\boldsymbol E^2
\end{equation}
in simple electrostatic units. A net charge leaves field extending far outside
the system, while a tightly bound neutral pair has only a short-range dipole
field. Neutralization therefore removes field energy from the exterior.

Imagine increasing the electromagnetic coupling enormously. The field-energy
penalty for net charge would become so large that low-energy objects would
overwhelmingly be neutral, and extracting a charged constituent could become
prohibitively expensive. Nambu transferred this intuition to color: the eight
color charges are sources of eight gluon fields, so a large color-field energy
could forbid isolated nonsinglets.

That intuition is essentially right, but QCD contains a decisive new feature.
The photon carries no electric charge, whereas the gluon carries color. The
photon is not a source of the classical electromagnetic field, and Maxwell's
equations are linear: two electromagnetic waves superpose and pass through
one another. A gluon transforms in the color octet, so gluons source and
interact with other gluons. Even different parts of a single gluon field
interact. The non-Abelian gauge field is therefore nonlinear.

\subsection{Flux Tubes and a Linear Potential}

At long distance the gluon self-interaction prevents color flux from spreading
out like an ordinary electric dipole field. The flux is squeezed into a tube
joining a quark to an anti-quark. The field strength and cross-sectional area
of a long tube are approximately fixed, so the tube stores an approximately
constant energy per unit length. If \(\sigma\) denotes this string tension,
then
\begin{equation}
V(R)\simeq\sigma R+\text{constant}
\qquad (R\ \text{large}).
\label{eq:l4-linear-potential}
\end{equation}

\begin{figure}[htbp]
\centering
\includegraphics[width=0.78\textwidth]{lecture_04_figure_06.png}
\caption{Color flux bundled into a long tube between a quark and an anti-quark.}
\lectureframe{01:31:35}
\end{figure}

This is stronger than the Coulomb potential, whose magnitude falls as \(1/R\).
Pulling the color sources farther apart does not weaken the force indefinitely;
it lengthens the tube and keeps adding energy. A lone color charge in an
infinite volume would require a tube extending without end and therefore
unbounded energy. That is the operational content of saying that only color
singlets have finite-energy asymptotic states.
\editorialnote{The string-tension notation and explicit long-distance potential make quantitative the source's statement that a uniform flux tube carries uniform energy per unit length.}

\subsection{String Breaking and Jets}

Suppose a hard collision strikes one quark inside a meson. The quark begins to
move away, but the flux tube stretches behind it. It may lose its kinetic
energy and be pulled back. More often, once the stored field energy is large
enough, the vacuum produces a new quark--anti-quark pair. The new anti-quark
terminates the flux attached to the escaping quark, while the new quark
terminates the flux attached to the original anti-quark. One stretched
singlet becomes two shorter singlets,
\begin{equation}
(q\bar q)_{\mathrm{stretched}}
\longrightarrow
(q\bar q)+(q\bar q).
\label{eq:l4-string-breaking}
\end{equation}

\begin{figure}[htbp]
\centering
\includegraphics[width=0.78\textwidth]{lecture_04_figure_07.png}
\caption{Pair creation breaks one stretched flux tube into two color-singlet tubes.}
\lectureframe{01:34:52}
\end{figure}

\begin{classroomqa}
\audiencequestion{If a struck quark has enough energy to keep stretching the tube, what else can happen besides its being pulled back?
\lecturetimestamp{01:33:50}}
\lecturerresponse{The tube can break by quark--anti-quark pair creation. Its energy becomes the mass and kinetic energy of the new pair. The endpoints reconnect into two singlets, so the process creates additional mesons rather than an isolated quark.}
\end{classroomqa}

If one of those partons still carries a large momentum, the process repeats.
The result is a cascade of color-singlet hadrons collimated around the original
parton direction: a jet. This conversion of energetic quarks and gluons into
hadrons is called hadronization. The initiating colored object is inferred
from the jet, never observed by itself.

\begin{classroomqa}
\audiencequestion{Does string breaking produce only mesons, or can other hadrons appear in the jet?
\lecturetimestamp{01:35:20}}
\lecturerresponse{Other hadrons can certainly appear. The only color requirement is that the final observable objects be singlets. Mesons are the simplest picture for one break, but baryons, antibaryons, and more complicated singlet combinations can also be produced.}
\end{classroomqa}

Three quarks can terminate their flux in a color-singlet junction. The
three-pronged picture is the dynamical counterpart of the antisymmetric
epsilon tensor in Equation~\eqref{eq:l4-baryon-singlet}: no net color flux
escapes from the complete baryon.

\begin{figure}[htbp]
\centering
\includegraphics[width=0.72\textwidth]{lecture_04_figure_08.png}
\caption{A three-pronged color-flux junction for a baryon, below the meson string-breaking sketch.}
\lectureframe{01:36:12}
\end{figure}

\section{Gauge Fields, Conserved Charges, and the Standard Model}

At the present level, a gauge theory may be recognized as a theory in which a
conserved charge sources a photon-like field. Electric charge sources the
electromagnetic field; the eight color charges source the eight gluon fields.
The deeper formulation makes the symmetry local and introduces a connection
that compares internal states at neighboring spacetime points, but that
construction is not needed for the confinement picture developed here.
\editorialnote{The local-symmetry formulation is the standard deeper meaning of gauge theory alluded to but deliberately postponed in the source.}

Symmetry, conserved charge, and flux are tightly connected. If field lines may
end only on charges, then a source cannot simply disappear while its flux
remains attached to nothing. Nor can all of its distant field disappear
instantaneously, because field changes propagate causally. In local form the
same physics is expressed by a continuity equation,
\begin{equation}
\frac{\partial\rho}{\partial t}
+\boldsymbol\nabla\cdot\boldsymbol j=0,
\end{equation}
which says that charge can leave a region only by flowing through its boundary.
The field-line argument is the geometric picture of that local conservation
law.

We can now locate the result inside the Standard Model. Color contributes the
eight generators of \(SU(3)\) and their eight gluon fields. The Abelian sector
provides one generator and one gauge field. Three more generators belong to an
additional \(SU(2)\), not to ordinary spatial angular momentum. Before symmetry
breaking the gauge structure is
\begin{equation}
SU(3)_c\times SU(2)_L\times U(1)_Y,
\end{equation}
with \(8+3+1=12\) generators. The next task is to understand the electroweak
\(SU(2)\), then the Higgs field and the way the observed photon, \(W^\pm\), and
\(Z\) emerge.
\editorialnote{The source counts the photon together with the eight gluons and describes the remaining factor as \(SU(2)\times U(1)\). In the unbroken Standard Model the Abelian factor is hypercharge \(U(1)_Y\); the photon appears only after electroweak mixing.}

The central chain is now complete. Abstract groups act through
representations; infinitesimal generators are the conserved charges; quarks,
anti-quarks, and gluons occupy the \(\boldsymbol 3\),
\(\bar{\boldsymbol 3}\), and \(\boldsymbol 8\); tensor products identify the
singlet hadrons; and the nonlinear gluon field turns singlet selection into
the physical phenomenon of confinement.
